% !TEX root = main.tex
Questo lavoro indaga se la componente appresa di un agente di Deep Reinforcement Learning aggiunga valore rispetto a una regola strutturale fissa che punta a regolare l'esposizione al mercato. 
L'agente addestrato con PPO su una rete ricorrente (GRU) e con politica basata su una distribuzione Beta, non sostituisce la strategia di volatility targeting di Moreira e Muir ma la corregge attraverso un unico moltiplicatore appreso $m_\theta$.
La selezione dei titoli è congelata a equal-weight in modo da isolare la sola decisione di esposizione e permettere un confronto pulito tra ciò che la rete impara e ciò che la formula fornisce.

Il modello è valutato out-of-sample su più regimi di mercato e su panieri diversi, con protocollo multi-seed e una frontiera di controllo a parità di liquidità media. 
I risultati mostrano che la parte appresa non aggiunge nulla che non si possa ottenere scegliendo i parametri strutturali fissi. 
Il minor drawdown dell'agente deriva dal mantenere più liquidità, non da un reale tempismo, l'agente quindi non ha svilupato alcun tipo d'intelligenza di market timing. 
L'agente ricade sulla frontiera della regola deterministica, resta dominato sullo Sharpe e la dipendenza dal regime dell'esponente non è sfruttabile in tempo reale.
La ragione è economica: la volatilità è prevedibile, la direzione del mercato no.

Il contributo non è dimostrare che l'apprendimento batte una regola, ma quantificare con rigore, tramite ablation e controlli, il limite dell'adattività appresa in un dominio a basso rapporto segnale/rumore. 
Come obiettivo secondario è stata sviluppata un'applicazione mobile local-first che esegue il modello on-device, rendendo le strategie studiate concretamente utilizzabili.
